\documentclass[12pt]{article}
\usepackage{multicol}
\usepackage[margin=1.25in]{geometry}
\setlength{\parskip}{0.5em}
\setlength{\columnsep}{0.4in}

\title{Reading Order Test: Two Columns Across Pages}
\author{legible-pdf synthesized corpus}
\date{}

\begin{document}
\maketitle

\section{Introduction}

This document is part of the legible-pdf synthesized corpus. Its purpose is
to exercise reading-order preservation across two-column layouts and across
page boundaries. Each numbered marker in the body below identifies a
specific sentence whose order in the document is known by construction from
the source. Converters that scramble reading order will produce markdown in
which these markers appear out of sequence.

\section{Body}

\begin{multicols}{2}

\noindent
\textbf{Marker 01.} The cardinal axiom of this experiment is that ordered
markers must remain ordered after conversion.

\textbf{Marker 02.} A converter that respects two-column reading order will
present marker two immediately after marker one.

\textbf{Marker 03.} Each marker introduces a distinct topical noun so that
probes can reference content unambiguously: this sentence mentions a
\emph{telescope}.

\textbf{Marker 04.} The fourth sentence mentions a \emph{kettle}, so that
probes can ask whether telescope precedes kettle.

\textbf{Marker 05.} A common failure mode is the converter treating each
column as a separate stream and reading column one of page one followed by
column one of page two.

\textbf{Marker 06.} A subtler failure mode is the converter weaving the two
columns together line-by-line, producing interleaved gibberish that nominally
contains all the words but in scrambled order.

\textbf{Marker 07.} Marker seven mentions a \emph{compass}, the third
distinct topical noun.

\textbf{Marker 08.} Marker eight begins the discussion of how layout
detection typically works in pipeline-based converters.

\textbf{Marker 09.} A layout detector identifies bounding boxes for text
regions and labels them as headings, paragraphs, columns, captions, or
figures.

\textbf{Marker 10.} A reading-order resolver then sorts those bounding boxes
into a sequence that approximates the order a human reader would follow.

\textbf{Marker 11.} Marker eleven mentions a \emph{lantern}, the fourth
distinct topical noun.

\textbf{Marker 12.} Resolution mistakes accumulate: missing one column
boundary on page one cascades into wrong order for the rest of the document.

\textbf{Marker 13.} The middle of the document is here. Markers below this
point are intended to span at least one page break in the rendered PDF.

\textbf{Marker 14.} Marker fourteen mentions a \emph{harpoon}, the fifth
distinct topical noun.

\textbf{Marker 15.} A correct reading-order resolver places marker fifteen
after marker fourteen, regardless of which page each occupies.

\textbf{Marker 16.} Page boundaries are a second axis of failure: even when
column order is right within a page, a converter may interleave pages
incorrectly.

\textbf{Marker 17.} Marker seventeen mentions a \emph{anchor}, the sixth
distinct topical noun.

\textbf{Marker 18.} The penultimate marker tests that the converter does not
truncate at page boundaries.

\textbf{Marker 19.} Marker nineteen mentions a \emph{sextant}, the seventh
distinct topical noun.

\textbf{Marker 20.} The final marker, marker twenty, must appear last in the
converted document. If it does not, the converter has failed reading-order
preservation in some way that the previous nineteen markers may help
diagnose.

\end{multicols}

\section{Conclusion}

The body above contains twenty sequential markers spanning two columns and
at least one page break. Markers also embed distinct topical nouns
(telescope, kettle, compass, lantern, harpoon, anchor, sextant) so that
probes may ask about content order in addition to marker order.

\end{document}
